\begin{abstract}
Static application security testing (SAST) tools remain essential for
detecting vulnerabilities early in the development lifecycle, yet
traditional analyzers suffer from high false-positive rates and shallow
semantic understanding. A new wave of LLM-based security scanners
promises deeper reasoning, but no open, reproducible benchmark exists to
measure their real-world effectiveness. We introduce
\textsc{RealVuln}, the first fully open-source vulnerability-detection
benchmark built from 26 real-world Python repositories containing 796
hand-labeled findings and 120 false-positive traps. We evaluate 15
scanners---3 traditional SAST tools, 11 LLM-based scanners, and 1
hybrid---using the recall-weighted F3 score ($\beta{=}3$, recall weighted 9$\times$ over precision) as the primary metric, with F2 ($\beta{=}2$) reported as a secondary metric for lower-risk contexts. Our
results show that the best LLM scanner (Claude Sonnet~4.6, F3\,=\,51.7)
outperforms the best traditional SAST tool (Semgrep, F3\,=\,17.7) by
nearly 3$\times$, with LLM-based tools achieving 2--4$\times$ higher F3
overall. Agentic multi-turn scanning adds approximately 46\% F2 over
single-turn inference for only ${\sim}$6\% additional cost. All code,
ground-truth data, scanner outputs, and scoring scripts are released
under an open-source license. \textsc{RealVuln} is designed as a living
benchmark---versioned, community-driven, and with a roadmap toward
multi-language coverage.
\end{abstract}
